%!TeX root = main.tex

The objective of this laboratory session is to elaborate a \texttt{Python} wrapper for the covariance matrix computation implemented in \texttt{C}. This enables to combine the strengths of both languages into a single workflow.

Scientific computing often faces a fundamental trade-off: high-level languages 
such as \texttt{Python} offer rapid development, readable syntax, and 
straightforward debugging, while low-level languages such as \texttt{C} provide 
direct memory control and computational efficiency. Rather than choosing one over the other, a common pattern in high performance computing (HPC) is to use each language for what it does best: \texttt{Python} as an orchestration layer, and \texttt{C} as the computational engine.

This approach is particularly relevant for covariance matrix computation over 
large datasets, where the mathematical operations are well-defined but 
computationally expensive. As demonstrated in previous reports, parallelisation 
strategies such as \texttt{OpenMP} allow \texttt{C} to exploit multi-core 
architectures efficiently. The role of \texttt{Python} here is not to replace 
this computation, but to provide a convenient interface to it.

To benchmark this approach, two implementations are compared: one using 
\texttt{Python}'s \texttt{NumPy} library, which provides vectorised operations 
on the full dataset, and one using the \texttt{C}/\texttt{OpenMP} function called directly from \texttt{Python} via a subprocess wrapper since, as demonstrated in previous reports, the OpenMP implementation achieved the best execution times among the parallelisation strategies evaluated. This comparison illustrates the performance gap between interpreted and compiled code, even when the interpreted solution uses optimised numerical libraries.

The wrapper is implemented as a \texttt{Python} class, \texttt{CovarianceOpenMP}, following an object-oriented design that encapsulates all communication with the \texttt{C} binary. The constructor stores the path to the executable and the desired number of threads, making the wrapper reusable across different configurations without modifying the calling code.

\begin{minted}[breaklines, fontsize=\small]{python3}
from __future__ import print_function, division
import subprocess
import time
import re
import os
import numpy as np

# C/MPI wrapper
class CovarianceOpenMP(object):
    def __init__(self, executable, n_threads = 4):
        self.executable = executable
        self.n_threads  = n_threads

    def compute(self):
        t_wall_start = time.time()
\end{minted}

The \texttt{compute()} method is responsible for launching the \texttt{C} 
process. Since \texttt{OpenMP} controls parallelism through environment variables, \texttt{OMP\_NUM\_THREADS} is set before the process starts. The binary is then launched using \texttt{subprocess.Popen}, with \texttt{stdout} and \texttt{stderr} redirected to pipes so that the output can be captured and parsed programmatically rather than printed directly to the terminal.

\begin{minted}[breaklines, fontsize=\small]{python3}
        env = os.environ.copy()
        env["OMP_NUM_THREADS"] = str(self.n_threads)

        proc = subprocess.Popen(
            [self.executable],
            stdout  = subprocess.PIPE,
            stderr  = subprocess.PIPE,
            env     = env
        )
\end{minted}

Once the process finishes, \texttt{proc.communicate()} blocks until all output 
is received and the process exits. The return code is then checked to detect 
failures early. Since \texttt{subprocess} returns raw bytes in some 
\texttt{Python} versions, the output is decoded to a \texttt{UTF-8} string 
before processing. Finally, two parsing methods extract the timing information 
and the covariance matrix values from the captured output, which are returned 
to the caller.

\begin{minted}[breaklines, fontsize=\small]{python3}
        print('Starting computation using the Python Wrapper with', self.n_threads, 'threads')
        stdout, stderr = proc.communicate()
        t_wall = time.time() - t_wall_start

        if proc.returncode != 0:
            raise RuntimeError(
                "Process failed (code {0}):\n{1}".format(proc.returncode, stderr)
            )

        if isinstance(stdout, bytes):
            stdout = stdout.decode("utf-8")

        timings = self._parse_timings(stdout)
        cov     = self._parse_cov(stdout)

        return cov, timings, t_wall
\end{minted}

Thanks to the \texttt{NumPy} \texttt{Python} library, the covariance matrix calculation is straightforward, including in its instructions, the \texttt{np.cov(matrix)} function that calculates the covariance matrix with only one command. Nonetheless, the \texttt{dwarf} node has an old installed version of \texttt{Python} (2.6). In this version, the \texttt{NumPy} library does not include the cubic square-root, so it needs to be implemented:

\begin{minted}[breaklines, fontsize=\small]{python3}
def cbrt(x):
    if x >= 0.0:
        return x ** (1.0 / 3.0)
    else:
        return -((-x) ** (1.0 / 3.0))
\end{minted}

With the already defined \texttt{cbrt(x)} function, the computed variables can be calculated.

\begin{minted}[breaklines, fontsize=\small]{python3}
def compute_derived(var):
    print("----- Computing derived variables -----")
    comp    = [None] * 5
    comp[0] = np.sin(var[0]) + np.cos(var[1])
    comp[1] = np.exp(var[2]) + np.exp(-1.0 * var[3])
    comp[2] = np.sin(var[0]) * np.cos(var[1]) + np.cos(var[2]) * np.sin(var[3])
    comp[3] = np.hypot(var[1], var[2])
    comp[4] = np.vectorize(cbrt)(var[4])
    print("----- Derived variables computed -----")
    return comp

# Python Covariance Calculation
def calc_cov(final_var, N):
    print("----- Computing covariance -----")
    matrix  = np.vstack(final_var)
    cov     = np.cov(matrix)
    print("----- Covariance computed -----")
    return cov
\end{minted}

The output of the wrapper is given in the following output:

\begin{minted}{bash}
Starting computation using the Python NumPy functions
Reading 134217728 samples...
----- Reading variables from files -----
----- Files read completed -----
  Read time: 3.7787 sec
----- Computing derived variables -----
----- Derived variables computed -----
  Compute variables time: 115.3501 sec
  Computing covariance...
----- Covariance computed -----
  Covariance time: 316.6341 sec
  cov(1,1): 0.333304

Starting computation using the Python Wrapper with 40 threads
  Wall time: 122.7729 sec
  Internal C timings:
    read  : 114.3998 sec
    comp  : 2.8731 sec
    mean  : 0.5191 sec
    cov   : 3.5948 sec
  cov[0,0]: 0.333304

------------------------------------------------------------
  SUMMARY (both on 134217728 samples)
------------------------------------------------------------
  NumPy  read : 3.7787 sec
  NumPy  comp : 115.3501 sec
  NumPy  cov  : 316.6341 sec
  C/OpenMP cov: 3.5948 sec
  Speedup (cov only): ~88.1x
============================================================
\end{minted}

The results reveal an interesting asymmetry between the two implementations. 
Regarding file reading, \texttt{NumPy}'s \texttt{np.fromfile} reads all 
134 million elements in 3.78~sec, significantly faster than the \texttt{C} 
sequential \texttt{fread} which required 114.40~sec. This is likely due to 
\texttt{NumPy} using optimised buffered I/O internally.

However, the advantage reverses decisively during computation. The derivation 
of the five computed variables, which takes only 2.87~sec in \texttt{C} with 
40 \texttt{OpenMP} threads, required 115.35~sec in \texttt{NumPy}, which can be explained by the overhead of applying \texttt{np.vectorize} to the 
scalar \texttt{cbrt} function, which negates \texttt{NumPy}'s vectorisation 
benefits. In contrast, when calculating the covariance matrix, \texttt{C}/\texttt{OpenMP} completes it in 3.59~sec, while \texttt{NumPy} requires 
316.63~sec, yielding a speedup of approximately $88\times$.

In conclusion, while \texttt{NumPy} provides competitive I/O performance, 
the \texttt{C}/\texttt{OpenMP} implementation is vastly superior for 
compute-intensive operations, demonstrating that compiled, parallelised code 
remains the appropriate choice for large-scale numerical computation.